\section{User Interface Design}
\label{sec:ui}

\subsection{Interaction Principles}

Four principles govern the interface, in priority order when they conflict.

\begin{enumerate}
\item \textbf{The image is the interface.} Controls overlay or dock; they never
crop the image below 55\% of the screen area. Any control that requires the
image to shrink is a modal sheet, not a permanent chrome element.
\item \textbf{One vocabulary.} The names in Table~\ref{tab:controlmap} appear
identically in the UI, in the tooltip, in the exported metadata, and in this
document. There is no simplified alias.
\item \textbf{Progressive disclosure without a beginner mode.} Every panel has
a primary row of two to four controls and a disclosure for the remainder. The
beginner and the professional use the same controls; the beginner sees fewer at
once.
\item \textbf{Every control teaches.} Long-pressing any control presents a
one-paragraph explanation of the physical process it corresponds to, with a
diagram where useful. This is the mechanism by which the product's central claim
--- that development is more truthful than filtering --- becomes something the
user experiences rather than reads in marketing copy.
\end{enumerate}

\subsection{Navigation Structure}

Three root workspaces in a tab structure: \textbf{Capture}, \textbf{Develop},
\textbf{Library}. Develop is reachable only with an active asset; selecting an
asset in Library or completing a capture pushes into it. This is a shallow
hierarchy by design --- the deepest navigation path in the application is three
levels (Library $\to$ Develop $\to$ Stage detail).

\begin{figure*}[htbp]
\centering
\begin{tikzpicture}[
  font=\tiny,
  phone/.style={draw=emLine, line width=0.5pt, rounded corners=3pt, minimum width=42mm, minimum height=78mm},
  ctl/.style={draw=emLine!60, fill=emGray, rounded corners=1pt, minimum height=3.2mm, inner sep=1pt, align=center},
  img/.style={draw=emLine!50, fill=emBlue!7, minimum width=38mm},
  lbl/.style={font=\tiny\bfseries},
  wfcap/.style={font=\scriptsize\bfseries, text=emLine}
]

%============ CAPTURE ============
\begin{scope}[shift={(0,0)}]
  \node[phone] (p1) at (0,0) {};
  \node[wfcap, above=1mm of p1.north] {CAPTURE};
  % status bar
  \node[ctl, minimum width=38mm, fill=white] at (0,3.62) {\tiny Portra 400 \; $\cdot$ \; 2383 \; $\cdot$ \; 135 \hfill ISO 400 \; 1/125 \; f/2.8};
  % viewfinder
  \node[img, minimum height=44mm] (vf) at (0,1.0) {};
  \node at (0,1.0) {\itshape live preview, recipe applied};
  % grid overlay hint
  \draw[emLine!25] (-1.9,-0.47) -- (1.9,-0.47);
  \draw[emLine!25] (-1.9,2.47) -- (1.9,2.47);
  \draw[emLine!25] (-0.63,-1.2) -- (-0.63,3.2);
  \draw[emLine!25] (0.63,-1.2) -- (0.63,3.2);
  % histogram
  \node[ctl, minimum width=17mm, minimum height=6mm, anchor=north west] at (-1.9,-1.35) {\tiny log histogram};
  \node[ctl, minimum width=19mm, minimum height=6mm, anchor=north east] at (1.9,-1.35) {\tiny zone / clip warn};
  % manual control strip
  \node[ctl, minimum width=8.5mm, anchor=north west] at (-1.9,-2.25) {ISO};
  \node[ctl, minimum width=8.5mm, anchor=north west] at (-0.95,-2.25) {SHUT};
  \node[ctl, minimum width=8.5mm, anchor=north west] at (0.0,-2.25) {WB};
  \node[ctl, minimum width=8.5mm, anchor=north west] at (0.95,-2.25) {FOCUS};
  % dial
  \draw[emLine!60] (0,-3.15) circle (0.55);
  \node at (0,-3.15) {\tiny stock};
  \node[ctl, minimum width=9mm, anchor=east] at (-0.8,-3.15) {recipes};
  \node[ctl, minimum width=9mm, anchor=west] at (0.8,-3.15) {last shot};
  % shutter
  \draw[emAcc, line width=0.8pt] (0,-3.15) circle (0.78);
\end{scope}

%============ DEVELOP ============
\begin{scope}[shift={(5.6,0)}]
  \node[phone] (p2) at (0,0) {};
  \node[wfcap, above=1mm of p2.north] {DEVELOP};
  \node[ctl, minimum width=38mm, fill=white] at (0,3.62) {\tiny $\leftarrow$ \hfill IMG\_4471 \hfill compare \; $\cdot$ \; export};
  \node[img, minimum height=38mm] (im) at (0,1.4) {};
  \node at (0,1.4) {\itshape rendered image};
  % stage rail
  \node[ctl, minimum width=38mm, minimum height=4mm] at (0,-0.85)
    {\tiny NEG $\cdot$ DEV $\cdot$ PRINT $\cdot$ COLOR $\cdot$ GRAIN $\cdot$ HAL $\cdot$ OPT $\cdot$ AGE};
  % active panel
  \node[draw=emLine!60, fill=emGray, rounded corners=2pt, minimum width=38mm, minimum height=22mm, anchor=north] (panel) at (0,-1.25) {};
  \node[lbl, anchor=north west] at (-1.85,-1.4) {PRINT};
  % printer lights
  \node[anchor=north west, align=left] at (-1.85,-1.75)
    {\tiny printer lights};
  \node[ctl, minimum width=11mm, anchor=north west] at (-1.85,-2.1) {R \; 25 \; $-|+$};
  \node[ctl, minimum width=11mm, anchor=north west] at (-0.6,-2.1) {G \; 25 \; $-|+$};
  \node[ctl, minimum width=11mm, anchor=north west] at (0.65,-2.1) {B \; 27 \; $-|+$};
  % density slider
  \node[anchor=north west] at (-1.85,-2.55) {\tiny print density};
  \draw[emLine] (-1.85,-2.95) -- (1.85,-2.95);
  \draw[emAcc, fill=emAcc] (0.25,-2.95) circle (0.09);
  \node[anchor=north east] at (1.85,-2.55) {\tiny $+3$ pts};
  % disclosure
  \node[ctl, minimum width=38mm, anchor=north] at (0,-3.15) {\tiny more: saturation density $\cdot$ roll-off $\cdot$ shadow lift $\cdot$ neutral axis \; $\vee$};
\end{scope}

%============ LIBRARY ============
\begin{scope}[shift={(11.2,0)}]
  \node[phone] (p3) at (0,0) {};
  \node[wfcap, above=1mm of p3.north] {LIBRARY};
  \node[ctl, minimum width=38mm, fill=white] at (0,3.62) {\tiny Projects \hfill select \; $\cdot$ \; $\cdots$};
  % segmented
  \node[ctl, minimum width=38mm] at (0,3.05) {\tiny ALL \; $\cdot$ \; ROLLS \; $\cdot$ \; RECIPES \; $\cdot$ \; STOCKS};
  % grid
  \foreach \r in {0,1,2,3} {
    \foreach \c in {0,1,2} {
      \node[img, minimum width=11.5mm, minimum height=11.5mm, anchor=north west]
        at (-1.85+\c*1.27, 2.7-\r*1.27) {};
    }
  }
  \node[anchor=north west] at (-1.8,2.65) {\tiny \textbf{Roll 04}};
  % footer
  \node[ctl, minimum width=38mm, anchor=south] at (0,-3.72) {\tiny 214 frames \; $\cdot$ \; 6 recipes \; $\cdot$ \; 3.1 GB};
\end{scope}

\end{tikzpicture}
\caption{Primary workspace wireframes. Capture keeps the manual control strip and the stock dial within thumb reach of the shutter. Develop devotes the upper 55\% to the image with a horizontal stage rail selecting the active panel below. Library is grid-first with roll grouping, since users organize by shooting session.}
\label{fig:wireframes}
\end{figure*}

\subsection{The Develop Workspace}

The stage rail in Figure~\ref{fig:wireframes} is the structural expression of
the pipeline. Selecting a stage swaps the panel below without changing the image
area, so comparisons across stages do not shift the frame.

Each panel follows an identical layout contract:

\begin{itemize}
\item A title matching the pipeline stage name in Figure~\ref{fig:pipeline}.
\item Two to four \emph{primary} controls, always visible.
\item A disclosure row listing the secondary controls by name, so their
existence is discoverable without being present.
\item A reset affordance scoped to the stage.
\item A bypass toggle, which disables the stage entirely (and, through dead
pass elimination in Section~\ref{sec:rendergraph}, removes its cost).
\end{itemize}

\begin{table}[htbp]
\caption{Panel Contents by Stage}
\label{tab:panels}
\begin{center}
\renewcommand{\arraystretch}{1.15}
\begin{tabularx}{\columnwidth}{@{}l>{\raggedright\arraybackslash}X@{}}
\toprule
\textbf{Panel} & \textbf{Primary / secondary} \\
\midrule
NEG & Stock, film speed rating, exposure comp. / format, layer balance readout, curve inspector \\
DEV & Push/pull, development time / temperature, agitation, chemistry, CI readout \\
PRINT & Printer lights R/G/B, print density / print stock, silver retention \\
COLOR & Saturation density, highlight roll-off / shadow lift, neutral axis warm, neutral axis tint \\
GRAIN & Amount, size / seed, clustering, chroma correlation \\
HAL & Intensity, threshold / radius, color bias, ring weight \\
OPT & Diffusion, vignetting / chromatic aberration, distortion, bloom, defocus \\
AGE & Expired age, dust / scratches, edge fog, light leak, streaking \\
\bottomrule
\end{tabularx}
\end{center}
\end{table}

\subsection{Control Idioms}

Three control types cover the entire parameter surface.

\paragraph{The stepper-slider} Used for all continuous parameters. A horizontal
track with a draggable thumb, a numeric readout in the parameter's physical unit,
and a detent at the physically neutral value. Dragging with one finger is coarse;
a second finger down engages a $5\times$ fine mode without lifting. Double-tap
resets to the stage default.

\paragraph{The printer-point stepper} Used only for printer lights. Integer
increment and decrement buttons flanking a value, with a hold-to-repeat. It is
deliberately \emph{not} a slider, because printer points are integers and
because the tactile discreteness reinforces the unit
(Section~\ref{sec:printerlights}).

\paragraph{The stock dial} A circular selector for negative and print stock,
showing a live thumbnail of the current image under each candidate as it passes
under the selection indicator. Rendering eight thumbnails at $128^2$ costs
approximately 1.2\,ms total, so the previews are live rather than cached ---
which is the single most persuasive interaction in the product, because the user
sees the stock's actual effect on their own image rather than on a sample.

\subsection{Comparison Modes}

FR-19 requires before/after. Three modes are provided:

\begin{itemize}
\item \textbf{Toggle.} Press and hold anywhere on the image to show the
unprocessed scene-referred image with only the output transform applied.
\item \textbf{Split.} A draggable vertical divider. The two sides render from
the same graph with different recipes, sharing all resources whose parameters
are unchanged --- so a split view of a printer-light change costs
$1.9\times$, not $2\times$, and a split of a grain change costs $1.2\times$.
\item \textbf{Pin to stage.} Compare the current state against the pipeline
output at any earlier stage. This is the pipeline inspector of FR-20 expressed
as a comparison rather than as a separate screen, and it is how a user comes to
understand what each stage contributes.
\end{itemize}

\subsection{The Pipeline Inspector}

A full-screen modal presenting the pipeline as the diagram of
Figure~\ref{fig:pipeline}, with a live thumbnail rendered at each node showing
the image state at that point. Tapping a node jumps to its panel; long-pressing
shows the governing equation and the current parameter values.

This screen is, deliberately, the most technically explicit surface in the
product. It is not the default view and most users will visit it rarely. Its
purpose is to make the claim in Section~\ref{sec:intro-thesis} inspectable: a
user who doubts that anything physical is happening can look.

\begin{figure}[htbp]
\centering
\begin{tikzpicture}[
  font=\tiny,
  nd/.style={draw=emLine!70, fill=emGray, rounded corners=1.5pt, minimum width=12mm, minimum height=9mm, align=center},
  th/.style={draw=emBlue!40, fill=emBlue!8, minimum width=11mm, minimum height=7mm},
  ar/.style={-{Latex[length=1.2mm]}, draw=emLine!70}
]
\foreach \i/\n [count=\j from 0] in {0/RAW, 1/HAL, 2/NEG, 3/DEV} {
  \node[th] (t\i) at (\j*1.55, 0.9) {};
  \node[nd, below=1mm of t\i] (n\i) {\n};
}
\foreach \i/\n [count=\j from 0] in {4/ILR, 5/GRN, 6/PRN, 7/OUT} {
  \node[th] (t\i) at (\j*1.55, -1.4) {};
  \node[nd, below=1mm of t\i] (n\i) {\n};
}
\foreach \a/\b in {n0/n1, n1/n2, n2/n3} \draw[ar] (\a) -- (\b);
\foreach \a/\b in {n4/n5, n5/n6, n6/n7} \draw[ar] (\a) -- (\b);
\draw[ar] (n3.east) -- ++(0.35,0) |- ($(t4.west)+(-0.35,-0.5)$) -- ++(0.35,0);
\node[anchor=north, align=center, text width=52mm] at (2.3,-2.55)
  {\scriptsize tap a node to open its panel \; $\cdot$ \; long-press for the model};
\end{tikzpicture}
\caption{Pipeline inspector layout. Each node shows the live image state at that point in the chain, rendered from the same graph with an early termination.}
\label{fig:inspector}
\end{figure}

\subsection{Gestures}

\begin{table}[htbp]
\caption{Gesture Map}
\label{tab:gestures}
\begin{center}
\renewcommand{\arraystretch}{1.15}
\begin{tabularx}{\columnwidth}{@{}l>{\raggedright\arraybackslash}X@{}}
\toprule
\textbf{Gesture} & \textbf{Action} \\
\midrule
Pinch & Zoom, 1$\times$ to 400\%. Triggers ROI rendering above 100\%. \\
Two-finger pan & Pan when zoomed. \\
Press and hold on image & Show original (toggle compare). \\
Double-tap image & Zoom to 100\% at tap point, or back to fit. \\
Horizontal swipe on stage rail & Change active panel. \\
Horizontal swipe on image & Previous/next asset (only when at fit zoom). \\
Two-finger drag on any slider & Fine adjustment, $5\times$ resolution. \\
Three-finger swipe left/right & Undo / redo (system convention). \\
Long-press control label & Explanation sheet. \\
\bottomrule
\end{tabularx}
\end{center}
\end{table}

Note the conflict resolution between horizontal image swipe (asset change) and
pan: asset swiping is enabled only at fit zoom, where panning is meaningless.
This is checked in the gesture recognizer's \texttt{shouldBegin}, not by
disabling recognizers, so the transition is seamless as the user zooms out.

\subsection{Visual Design}

The visual language is quiet, achromatic, and deferential to the image. The
chrome is a single neutral gray at two elevations, with one accent color used
exclusively for modified-from-default state. This is not minimalism for its own
sake: any chroma in the surrounding UI biases the perception of the image being
graded, which is why professional grading suites are painted neutral gray. The
same reasoning that governs a color suite governs this interface, and it is
stated in the design system documentation so that it survives future
redesigns.

Explicitly excluded: leather textures, film-canister iconography, simulated
light leaks in the chrome, faux-analog typography, sprocket-hole borders, and
any skeuomorphic camera imagery. The application is a laboratory instrument, and
laboratory instruments do not cosplay.

\subsection{Dark and Light Appearance}

Both are supported. Dark is the default, because a dark surround improves the
perception of highlight roll-off and because the application is most often used
in editing rather than reference-viewing contexts. The image area's immediate
surround is held at a fixed 18\% neutral in \emph{both} appearances, since a
pure-black or pure-white surround measurably shifts perceived contrast; only the
outer chrome changes with appearance.

\subsection{Accessibility}

NFR-11 is met by: VoiceOver labels for every control stating name, value, and
unit ("printer light red, 25 points"); custom rotors for stage navigation and
for parameter adjustment; Dynamic Type support to XXL in all text, with panels
reflowing to a vertical layout at large sizes; Reduce Motion honored by
replacing panel transitions with cross-fades; and a minimum 44$\times$44 point
touch target for every interactive element including the printer-point steppers.

Color is never the sole carrier of information. The modified-from-default
indicator is an accent color \emph{and} a shape change on the slider detent.
Clipping warnings use a diagonal hatch pattern rather than a color overlay.
